\section{Experimental setup}
\label{sec:method}

\textbf{Platform.} We use one NVIDIA GeForce RTX 5090 (Blackwell, 170 SMs, 32\,GB GDDR7, 575\,W) in a
workstation with an Intel Core Ultra~9 285K and 64\,GB of RAM, running Ubuntu~24.04, PyTorch~2.13
(CUDA~13.0), Transformers~5.19, bitsandbytes~0.50.2 and vLLM~0.32 (CUDA graphs on). Our
microbenchmarks measure \BWmeas\,TB/s of memory bandwidth and \BFmeas{} (BF16) and \FPeightmeas{}
(FP8) TFLOP/s for large matrix multiplies, which puts the BF16 ridge point at about \RidgeBF{}
FLOP/byte.

\textbf{Models and configurations.} The main study uses instruction-tuned
Qwen2.5~\cite{qwen2024qwen25} at 3B, 7B and 14B parameters; Qwen3-8B and Qwen3-32B~\cite{qwen3} test
generality. \Cref{tab:configs} lists the eight configurations. The Hugging Face (HF) backend is what a
user gets with one extra argument to \texttt{from\_pretrained}: bitsandbytes quantizes the BF16
checkpoint while loading, and PyTorch runs the model eagerly. The vLLM backend serves the official
Qwen GPTQ and AWQ checkpoints with Marlin kernels and converts BF16 weights to FP8 while loading. The
embedding, the output layer (LM head) and the norms stay in 16 bits in every configuration.

\begin{table}[t]
\centering\small
\caption{The eight configurations. W/A: bits of the weights and of the activations in the matrix
multiply.}
\label{tab:configs}
\setlength{\tabcolsep}{4.5pt}
\begin{tabular}{llllll}
\toprule
Backend & Name & W/A & Method & Kernel & Weights from \\
\midrule
HF   & BF16      & 16/16           & ---            & cuBLAS                 & BF16 checkpoint \\
HF   & INT8      & 8/8$^{\dagger}$ & LLM.int8()     & bitsandbytes, cuBLASLt & quantized on load \\
HF   & NF4       & 4/16            & NF4, block 64  & bitsandbytes           & quantized on load \\
vLLM & BF16      & 16/16           & ---            & cuBLAS                 & BF16 checkpoint \\
vLLM & FP8       & 8/8             & FP8 E4M3       & CUTLASS                & quantized on load \\
vLLM & GPTQ-INT8 & 8/16            & GPTQ, group 128 & Marlin                & Qwen checkpoint \\
vLLM & AWQ-INT4  & 4/16            & AWQ, group 128  & Marlin                & Qwen checkpoint \\
vLLM & GPTQ-INT4 & 4/16            & GPTQ, group 128 & Marlin                & Qwen checkpoint \\
\bottomrule
\end{tabular}
\par\vspace{3pt}\parbox{\linewidth}{\scriptsize $^{\dagger}$Activation columns with values above 6.0
are multiplied in FP16.}
\end{table}

\textbf{Workloads.} \emph{Controlled sweep:} each configuration generates exactly 128 tokens for
$B\in\{1,2,4,\dots,64\}$ prompts (up to 256 on vLLM) of $L\in\{128,512,1024,2048\}$ random tokens,
so every configuration does the same work; a smaller sweep varies the output length. \emph{Real traffic:} 300 ShareGPT~\cite{sharegpt} conversations (median prompt 31 tokens,
median answer 284 tokens), served by vLLM to 1--64 concurrent clients and by HF in static batches of
16. \emph{Quality:} WikiText-2~\cite{merity2017wikitext} perplexity on 2048-token windows;
GSM8K~\cite{cobbe2021gsm8k} on all 1,319 test questions with zero-shot chain of thought and greedy
decoding; MMLU~\cite{hendrycks2021mmlu} on all 14,042 test questions, zero-shot. Prompts use each
model's chat template, and Qwen3 runs in non-thinking mode.

\textbf{Metrics and statistics.} TTFT is the time to the first output token, TPOT the mean time of
each later token, and throughput the number of output tokens per second. Each cell has one warm-up and
three measured runs with fresh prompts; the run-to-run variation of the end-to-end time is below
\CVmax\% in 95\% of the cells and \CVabsmax\% at most. Memory and energy come from PyTorch and NVML. For quality, all configurations answer the same
questions, so we compare each with HF BF16 by an exact McNemar test, with a Bonferroni correction over
the \MNumTests{} comparisons of each benchmark ($\alpha=\BonfAlpha$).
